\section{Persona Panel, Prompt, and Reproducibility Details}
\label{sec:appendix}

\paragraph{Persona panel and how it was constructed.}
The panel is $P{=}10$ fixed personas. They were authored from aggregate demographics by a
deterministic, documented procedure rather than ad hoc, to limit authorial bias:
(i) we took the audience descriptors reported by Pew for U.S.\ social-media news consumers in the
Upworthy era~\citep{holcomb2013newsuse}---age band, gender, primary platform (Facebook, Twitter,
YouTube, Reddit), and news-engagement level (incidental/moderate/active);
(ii) we allocated the ten persona slots to cover this joint space in rough proportion to the
reported population (e.g.\ Facebook as the dominant news gateway, Twitter news consumers skewing
younger, Facebook news consumers skewing female);
(iii) for each cell we wrote a single behavioural sentence describing \emph{clicking} tendencies,
deliberately avoiding brand-specific content; the panel is authored from public demographics and
is not reused from any product. The complete, versioned panel ships in the replication package
(\texttt{data/persona\_panel.json}) with the grounding facts recorded inline; the procedure is
fully specified there so the panel can be regenerated or extended. We treat persona
\emph{psychographics} (beyond demographics) as out of scope and a direction for future work.

\paragraph{Elicitation prompt (verbatim template).}
\begin{quote}\small
You are simulating a specific person reacting to a news headline on social media in 2014.\\
Person: \{persona description\} (age \{band\}, \{gender\}, mainly on \{platform\}, news
engagement: \{level\}).\\
Headline: ``\{headline\}''\\
As this exact person, how likely are you to click this headline? Respond ONLY with JSON:
\{"click\_intent": <number between 0 and 1>\}.
\end{quote}

\paragraph{Predicted score.}
$\texttt{pred\_score}(v)=\frac{1}{P D}\sum_{i,d}\texttt{click\_intent}_{i,d}(v)$ over $P$
personas and $D$ draws. Variants are ranked by \texttt{pred\_score} within each package.

\paragraph{Seeds.}
Corpus sampling seed $42$; bootstrap seed $42$ ($10{,}000$ resamples); LLM generation seed
$42{+}d$ for draw $d$ (forwarded to the model's \texttt{generationConfig}). Each prediction
record stores model id, seed, temperature ($0.8$), draw count, and panel size.

\paragraph{Coverage and parse failures.}
LLM JSON parse failures are rare ($<1\%$ of calls) and retried; a package enters the metrics only
when \emph{all} its variants are scored (matching the primary join), so no package is ever ranked
from partial coverage. Per-persona draws that fail are simply averaged over the remaining draws.

\paragraph{Ground-truth reliability test.}
A package has a reliable winner when its top-CTR headline beats the runner-up by a one-sided
two-proportion $z$-test at $p<0.05$~\citep{kohavi2009controlled}. \reliableFraction{} of eligible
A/B tests pass this filter.

\paragraph{Model comparison (pilot).}
Table~\ref{tab:model_comparison} reports the pilot that motivated the model choice: three
\texttt{gemini} tiers on a common subset, with overlapping validity intervals.
% AUTO-GENERATED by scripts/generate_pilot_table.py — DO NOT EDIT BY HAND.
\begin{table*}[t]
  \centering
  \caption{Pilot model comparison on a common subset ($n=100$ packages). The validity intervals overlap across model tiers, so the cheapest model was selected. 95\% bootstrap CIs.}
  \label{tab:model_comparison}
  \small
  \begin{tabular}{lccc}
    \toprule
    Model & Top-1 & Kendall $\tau$ & Spearman $\rho$ \\
    \midrule
    \texttt{gemini-2.5-flash} & 0.390 [0.29, 0.49] & 0.116 [-0.01, 0.24] & 0.133 [-0.01, 0.27] \\
    \texttt{gemini-3.1-flash-lite} & 0.330 [0.24, 0.42] & 0.035 [-0.10, 0.17] & 0.051 [-0.09, 0.19] \\
    \texttt{gemini-3.5-flash} & 0.390 [0.29, 0.49] & 0.122 [-0.01, 0.25] & 0.139 [-0.00, 0.28] \\
    \bottomrule
  \end{tabular}
\end{table*}


\paragraph{Cross-model robustness.}
Table~\ref{tab:cross_model_gap} shows the persona-vs-baseline gap on all three model tiers; the
baseline wins on every one. All runs disable model thinking (\texttt{thinkingBudget}=0): the task
is a single scalar rating, and \texttt{flash-lite} performs no thinking by default, so this makes
the configuration explicit and the comparison fair and inexpensive.
% AUTO-GENERATED by scripts/generate_crossmodel_table.py — DO NOT EDIT BY HAND.
\begin{table*}[t]
  \centering
  \caption{Cross-model robustness ($n=100$ reliable packages, Kendall $\tau$). The no-persona baseline beats the persona panel on every model tier; persona conditioning hurts regardless of model.}
  \label{tab:cross_model_gap}
  \small
  \begin{tabular}{lccc}
    \toprule
    Model & Persona $\tau$ & Baseline $\tau$ & Gap \\
    \midrule
    \texttt{gemini-3.1-flash-lite} & 0.035 & 0.401 & $+0.367$ \\
    \texttt{gemini-2.5-flash} & 0.116 & 0.219 & $+0.103$ \\
    \texttt{gemini-3.5-flash} & 0.122 & 0.404 & $+0.282$ \\
    \bottomrule
  \end{tabular}
\end{table*}


\paragraph{Error-analysis examples.}
Table~\ref{tab:error_examples} shows real winning headlines the no-persona baseline ranked first
but the panel did not, illustrating the systematic mis-ordering discussed in
Section~\ref{sec:discussion}.
% AUTO-GENERATED by scripts/generate_supplementary_tables.py — DO NOT EDIT BY HAND.
\begin{table}[t]
  \centering
  \caption{Examples of real winning headlines that the no-persona baseline ranked first but the persona panel did not (of \nBaseRightPersonaWrong{} such cases). These are typically curiosity-gap / feel-good headlines the base model recognises as clickable.}
  \label{tab:error_examples}
  \small
  \begin{tabular}{p{5.6cm}cc}
    \toprule
    Real winning headline & CTR & \#var \\
    \midrule
    Bookmark This For When You’re Sad And Instantly Be 100 Times Happier & 1.38\% & 4 \\
    10 Amazing International Photos That Will Give You Some Perspective & 1.92\% & 4 \\
    If You Saw Someone Bullying A Kid For Being Gay Would You Speak Out... & 1.11\% & 2 \\
    If You Watch Hilary Swank Annoy A Whiny Senator, \$1 Will Go To Pun... & 1.67\% & 3 \\
    Seven Wonderful People Share Secrets That Are Killing Them.  This C... & 1.99\% & 4 \\
    A 9-Year-Old Asks McDonald’s CEO A Simple Question. She’s Not Lovin... & 3.23\% & 4 \\
    Did These Words Really Come Out Of Conan's Mouth? & 4.45\% & 3 \\
    We Use This Word So Often We Probably Don't Realize We're Saying It... & 1.68\% & 4 \\
    Did You Know You Can Go For A Swim In Google Earth? & 0.88\% & 3 \\
    Cookie Monster Is In Jail? I'm Actually Really Happy About The Reas... & 2.30\% & 4 \\
    \bottomrule
  \end{tabular}
\end{table}


\paragraph{Full persona panel.}
Table~\ref{tab:persona_panel} lists the ten personas; full behavioural descriptions ship in the
replication package (\texttt{data/persona\_panel.json}).
% AUTO-GENERATED by scripts/generate_supplementary_tables.py — DO NOT EDIT BY HAND.
\begin{table*}[t]
  \centering
  \caption{The full fixed persona panel ($P{=}10$), grounded in the Upworthy-era audience demographics~\citep{holcomb2013newsuse}, with the verbatim behavioural descriptions used in the prompt (for replicability). News engagement and political lean are in the replication package (\texttt{data/persona\_panel.json}).}
  \label{tab:persona_panel}
  \small
  \begin{tabular}{lp{2.6cm}p{10.5cm}}
    \toprule
    ID & Demographics & Behavioural description (verbatim) \\
    \midrule
    p01 & 18-29, F, Facebook & 29-year-old woman, scrolls Facebook on her phone during breaks; clicks human-interest and feel-good stories shared by friends; rarely seeks news directly. \\
    p02 & 18-29, M, Twitter & 24-year-old man, heavy Twitter user, follows journalists and culture accounts; clicks surprising or contrarian takes; low tolerance for clickbait he's seen before. \\
    p03 & 30-49, F, Facebook & 38-year-old working mother, checks Facebook in the evening; drawn to headlines about family, health, and local impact; shares uplifting or outrage-inducing stories. \\
    p04 & 30-49, M, Reddit & 41-year-old man, browses Reddit and Facebook; skeptical of sensational headlines, clicks when a headline promises concrete information or a strong narrative payoff. \\
    p05 & 50-64, F, Facebook & 57-year-old woman, mostly on Facebook; clicks emotionally resonant stories and anything about social causes she follows; less drawn to internet-slang headlines. \\
    p06 & 50-64, M, Facebook & 60-year-old man, reads news on Facebook and cable; clicks political and economic headlines, distrusts overtly hyped 'you won't believe' framing. \\
    p07 & 18-29, F, YouTube & 22-year-old college student, video-first, on YouTube and Instagram; clicks visually evocative or identity-relevant headlines; very sensitive to authenticity. \\
    p08 & 30-49, M, Facebook & 45-year-old man, casual Facebook user; clicks curiosity-gap headlines that promise a quick interesting fact, especially maps, lists, and 'what happened next' stories. \\
    p09 & 65+, F, Facebook & 68-year-old retiree, active on Facebook with family; clicks heartwarming, nostalgic, or cause-driven stories; avoids headlines that feel like tricks. \\
    p10 & 30-49, F, Twitter & 34-year-old professional, follows news and advocacy on Twitter; clicks data-driven or justice-oriented angles; fatigued by repetitive viral formats. \\
    \bottomrule
  \end{tabular}
\end{table*}

